%% EXHALE: Initial-condition benchmark (hydrostatic vs Parker-wind starts)
%% Author: Kwang-Il Seon
\documentclass[english,a4paper]{my_memo}
\synctex=1
\usepackage{listings}
\usepackage{xcolor}
\usepackage{booktabs}
\usepackage{amsmath}
\usepackage{graphicx}
\usepackage{babel}
\emergencystretch=3em
\lstset{basicstyle=\ttfamily\small, commentstyle=\color{gray}\itshape,
  frame=single, breaklines=true, columns=fullflexible, showstringspaces=false}

\begin{document}

\title{EXHALE: Does starting from a Parker-wind initial condition\\
       speed up (and stabilize) convergence?}
\author{Kwang-Il Seon}
\date{2026-06-07}
\maketitle

\begin{quote}\small\itshape
We test whether seeding the time integration with a transonic (Parker-type) wind
profile --- in particular a \emph{hot} Parker wind that also carries an elevated
temperature and ionization, as a \texttt{p-winds}-style initial guess --- reduces
the number of iterations (and wall-clock) needed to reach the converged
photoevaporative wind, compared with the default cold hydrostatic initial
condition (IC). A new \texttt{Hot Parker IC} option was added to
\texttt{set\_IC.f90}; the benchmark runs the same planet (HD\,209458b, H/He,
spherical) from three ICs and compares convergence and the IC-vs-final profiles.
\textbf{Bottom line:} the naive full-Parker-wind seed \emph{diverges} (its isothermal
density is globally inconsistent with the non-isothermal equilibrium and drains onto the
base; \S5--6); a corrected warm-seed IC that keeps the stable hydrostatic density and warms
only the thermal/ionization/velocity overlay \emph{converges to the identical steady state}
(\S7), but gives \emph{no speedup} for this planet, because ATES's convergence is gated by
hydrodynamic momentum-steadiness, not the thermal relaxation the seed shortcuts. The cold
hydrostatic IC remains the recommended default.

\medskip
\textbf{Naming.} This memo predates the rename: ``ATES'' below means this code,
now EXHALE, except where it explicitly cites the upstream ATES of Caldiroli et
al.\ (2021). The diagnostic hook is now \texttt{EXHALE\_DIAG\_BASE=1}.
\end{quote}

\tableofcontents
\bigskip

% =====================================================================
\section{Why the initial condition matters}
% =====================================================================
ATES integrates the 1-D radiation--hydrodynamics to a \emph{steady state}; the
initial condition only sets where the relaxation starts, not the final answer. But
it sets \emph{how long} the relaxation takes. The cost has two parts:
\begin{enumerate}\itemsep2pt
\item \textbf{Hydrodynamic relaxation} --- the velocity and density settle into the
  transonic outflow topology (subsonic base $\to$ sonic point $\to$ supersonic
  wind, with $\rho\,v\,r^2=$ const).
\item \textbf{Thermal/ionization relaxation} --- the cold, nearly-neutral gas is
  photo-heated to $\sim10^4$\,K and ionized; the temperature and ionization
  profiles settle to radiative/advective balance. This is the \emph{stiff} part,
  and historically the dominant cost (it motivated the semi-implicit energy solver;
  see \texttt{Update\_EXHALE\_solver}).
\end{enumerate}

The current default IC (\texttt{set\_IC.f90}) is an \textbf{isothermal hydrostatic
atmosphere}: $\rho(r)=\rho_{\rm bc}\,e^{b_{0,\rm eff}(-\phi(r)+\phi(0))}$, a near-zero
velocity seed $v=\tfrac12(r-r_0)$, $T=T_0$ (cold base temperature), and gas that is
\emph{mostly neutral}. The wind must therefore launch from rest \emph{and} heat/ionize
from cold/neutral --- both relaxations run from far away.

% =====================================================================
\section{Parker / \texttt{p-winds} as an initial guess}
% =====================================================================
A natural idea is to start from an analytic transonic wind. The isothermal
\textbf{Parker wind} is the algebraic solution of the steady isothermal momentum
equation in the gravitational (or Roche) potential,
\begin{equation}
  \tfrac12 v^2 - c^2\ln v - 2c^2\ln r + \phi(r) = B ,
  \label{eq:bernoulli}
\end{equation}
with $c^2$ the (dimensionless) isothermal sound speed$^2$ and $B$ fixed by the
sonic point $r_c$ where $\mathrm d\phi/\mathrm dr = 2c^2/r_c$; the wind takes the
subsonic root for $r<r_c$ and the supersonic root for $r>r_c$, and
$\rho=\dot M/(v\,r^2)$.

\texttt{p-winds} (dos Santos et al.\ 2022) is a popular open-source implementation
of exactly this: an isothermal Parker wind plus a steady-state H/He ionization
balance (and the He\,$2^3$S triplet). Crucially, \texttt{p-winds} \emph{prescribes}
the temperature $T$ and mass-loss rate $\dot M$ (free parameters fit to data); it
does \emph{not} solve the energy equation. Used as an ATES IC it would supply:
\begin{itemize}\itemsep1pt
\item the transonic $v(r)$, $\rho(r)$ topology (shortcuts hydrodynamic relaxation);
\item \textbf{and the ionization structure} (shortcuts part of the ionization
  relaxation) --- the key advantage over a bare cold Parker wind;
\item at a \textbf{warm} temperature (right velocity scale in the heated region),
\end{itemize}
but it cannot supply the \emph{non-isothermal} $T(r)$ of the real wind (cold base
$T_0\to\sim10^4$\,K plateau), and its $T,\dot M$ must be chosen.

\paragraph{Design choice.} EXHALE already contains an isothermal Parker-wind solver
(the \texttt{Transonic IC} option: Eq.~\ref{eq:bernoulli} at the cold $T_0$ sound
speed). Its limitation is that it is cold (wrong velocity scale) and neutral. We
therefore capture the \texttt{p-winds} advantage \emph{without an external
dependency} by generalizing it to a \textbf{hot Parker IC}: the same Parker solve at
a warm effective sound speed, plus a temperature and ionization ramp.

% =====================================================================
\section{The hot-Parker IC in EXHALE}
% =====================================================================
Implemented in \texttt{src/modules/init/set\_IC.f90}, enabled by an optional
\texttt{input.inp} line
\begin{lstlisting}
Hot Parker IC: 10000        # effective wind temperature T_wind [K]
\end{lstlisting}
(\texttt{hot\_parker\_ic}, \texttt{T\_wind\_ic} in \texttt{parameters.f90}; parsed
in \texttt{input\_read.f90}; auto-enables \texttt{Force start}).

\paragraph{Velocity / density.} The Parker solver \texttt{wind\_profile} (now taking
$c^2$ as an argument) is called with a \emph{warm} sound speed
\begin{equation}
  c^2_{\rm hot} = \frac{1+\delta p_{\rm bc}}{\rho_{\rm bc}}\;\frac{T_{\rm wind}}{T_0}\;\times 2 ,
\end{equation}
where the factor $2$ approximates the mean-molecular-weight halving when H is
ionized ($\mu\!\to\!\mu/2$). $v(r)$ is the transonic root of
Eq.~\ref{eq:bernoulli} and $\rho(r)=\dot M/(v r^2)$ pinned to $\rho_{\rm bc}$ at the
base. If the warm sonic point falls outside the domain the routine returns failure
and \texttt{set\_IC} falls back to the hydrostatic IC.

\paragraph{Temperature / ionization ramp.} With $\xi(r)=\min(v/c_{\rm hot},1)$ (0 at
the subsonic base $\to$ 1 in the supersonic wind),
\begin{align}
  T(r) &= T_0\,[\,1 + (T_{\rm wind}/T_0 - 1)\,\xi\,], \\
  x_{\rm HII}(r) &= \delta p_{\rm bc} + (1-\delta p_{\rm bc})\,\xi, \qquad
  x_{\rm HeII}(r) = \xi ,
\end{align}
and the pressure is kept thermodynamically consistent,
$p = \rho\,c^2_{\rm base}\,(T/T_0)\,(1+\xi)$, so that $p/\rho\to c^2_{\rm hot}$ in the
supersonic wind and $\to c^2_{\rm base}$ at the cold base (matching the base BC).
Metals are left neutral (they re-equilibrate per cell in one step). The base
boundary condition ($T=T_0$, $\rho_{\rm bc}$) is unchanged, so a thin
base-reconciliation layer still relaxes.

\paragraph{IC dump.} \texttt{set\_IC} now writes the (dimensional) initial profile
to \texttt{output/IC\_dump.txt} ($r$, $n$, $v$, $p$, $T$) for any IC, so the IC can
be plotted directly against the converged result.

% =====================================================================
\section{Benchmark setup}
% =====================================================================
\textbf{Planet:} HD\,209458b-like hot Jupiter (the standard ATES test): $R_p=1.38\,R_J$,
$M_p=0.69\,M_J$, $T_0=1320$\,K, $a=0.047$\,AU, $M_\star=1.13\,M_\odot$, power-law
EUV+X-ray SED ($\log_{10}L_{\rm EUV}=29.3$), \textbf{H/He only} (no metals, for a
clean and fast comparison), \textbf{spherical} domain to $10\,R_p$. Solvers: the
current defaults (semi-implicit energy, Brent $T$, analytic-Jacobian Newton).

\textbf{Three initial conditions} (same planet, same numerics):
\begin{enumerate}\itemsep1pt
\item \texttt{hydrostatic} --- the default cold hydrostatic IC.
\item \texttt{transonic}  --- the cold ($T_0$) Parker IC (\texttt{Transonic IC: True}).
\item \texttt{hot\_parker} --- the warm Parker IC (\texttt{Hot Parker IC: 10000}).
\end{enumerate}

\textbf{Measured:} iterations to convergence (the \texttt{count} at which $du$/$dtu$/
stall stops the run), wall-clock (\texttt{OMP\_NUM\_THREADS=4}), and the IC-vs-final
profiles ($T$, $v$, $\rho$ vs $r$). Files live in \texttt{ic\_start\_benchmark/}
(one subdirectory per IC); the figures are produced by
\texttt{ic\_start\_benchmark/ic\_compare.ipynb}.

% =====================================================================
\section{Results}
% =====================================================================
The three runs are summarized in Table~\ref{tab:icbench}. The headline result is
\emph{negative and instructive}: neither Parker-wind start helped, and the warm
(hot-Parker) start actively destabilized the integration.

\begin{table}[ht]
\centering
\caption{Three initial conditions, same planet (HD\,209458b, H/He, spherical, $10\,R_p$),
identical numerics (semi-implicit energy, Brent $T$, analytic-Jacobian Newton,
\texttt{OMP\_NUM\_THREADS=4}).}
\label{tab:icbench}
\small
\begin{tabular}{@{}llrrl@{}}
\toprule
IC requested & IC actually used & Steps & Wall [s] & Outcome \\
\midrule
\texttt{hydrostatic} & hydrostatic (cold) & 95\,274 & 1828.8 & converged ($du{=}2.0\mathrm{e}{-}2$, $dtu{=}6.2\mathrm{e}{-}7$) \\
\texttt{transonic}   & hydrostatic (fallback) & 95\,274 & 1836.7 & identical to hydrostatic \\
\texttt{hot\_parker} & warm Parker, $T_{\rm wind}{=}10^4$\,K & 30\,421$^\ast$ & 604.8$^\ast$ & \textbf{diverged} --- NaN at base \\
\bottomrule
\end{tabular}

\smallskip
{\footnotesize $^\ast$ The hot-Parker step count and wall-clock reflect an \emph{early
numerical breakdown}, \emph{not} convergence: the run terminated on a NaN-corrupted
state ($du{=}5.1\mathrm{e}{+}2$, $dtu{=}1.3\mathrm{e}{-}2$; $T{>}10^{40}$\,K at $r\!\approx\!1.0002$).
It must not be read as a speedup.}
\end{table}

\paragraph{transonic $\equiv$ hydrostatic.} At the cold base sound speed
$c_0^2=kT_0/\mu$, the Parker sonic point $r_c$ (where
$\mathrm d\phi/\mathrm dr=2c_0^2/r_c$) falls \emph{outside} the computational domain
for HD\,209458b --- there is no interior transonic root. \texttt{set\_IC} detects this
(``Transonic IC requested but no interior sonic point; falling back to hydrostatic
IC'') and uses the hydrostatic IC, so the two runs are byte-identical. A \emph{cold}
Parker IC offers nothing here.

\paragraph{hot\_parker diverges at the base.} The warm-Parker IC itself is smooth and
physical (Fig.~\ref{fig:icfinal}, dashed blue): a subsonic warm base
($v\!\approx\!0.7$\,km\,s$^{-1}$, $T\!\approx\!1840$\,K, $n\!\approx\!6\times10^{12}$)
ramping monotonically to a supersonic $10^4$\,K wind
($v\!\approx\!23$\,km\,s$^{-1}$ at $10\,R_p$), with no NaN in the IC dump. But during
time integration the base \emph{cannot reconcile} the warm, outflowing interior with
the cold fixed base boundary condition ($T=T_0=1320$\,K, $\rho=\rho_{\rm bc}$): within a
few cells of $r=1$ the base develops an \emph{inflow} (``breathing'', final
$v\!\approx\!-1.1\times10^5$\,cm\,s$^{-1}$), the cell-by-cell energy solve latches onto the
\emph{spurious hot root}, and $T$ runs away to $>10^{40}$\,K. This is the exact
breathing-base hot-root pathology documented for the \texttt{\_adv} post-processor (the
option-(c) guard; see \texttt{Update\_EXHALE\_solver}) --- but here it strikes the
\emph{main} time integration, where no such guard exists.

% =====================================================================
\begin{figure}[ht]
\centering
\includegraphics[width=\linewidth]{../benchmarks/ic_start_benchmark/ic_vs_final.pdf}
\caption{Initial condition (dashed) vs.\ converged final (solid) for each IC:
temperature (top), velocity (middle), density (bottom). \textbf{hydrostatic} and
\textbf{transonic} (the latter identical, having fallen back) relax from the flat cold
IC to the peaked $\sim10^4$\,K transonic wind. \textbf{hot\_parker} starts from a sane
warm-wind IC but its final state diverges (note the $10^{134}$\,K temperature axis and
the ``NaN at base'' annotation; the final velocity turns negative --- base inflow).}
\label{fig:icfinal}
\end{figure}

\begin{figure}[ht]
\centering
\includegraphics[width=0.92\linewidth]{../benchmarks/ic_start_benchmark/finals_overlay.pdf}
\caption{Converged final temperature and velocity for the ICs that converged
(hydrostatic and transonic): they coincide, confirming the IC affects only convergence
cost, not the steady state. The diverged hot-Parker run is excluded.}
\label{fig:finals}
\end{figure}

\begin{figure}[ht]
\centering
\includegraphics[width=0.92\linewidth]{../benchmarks/ic_start_benchmark/convergence_bars.pdf}
\caption{Steps and wall-clock per IC. The hot-Parker bars (red) are an early numerical
breakdown, \emph{not} a speedup --- the run never converged.}
\label{fig:bars}
\end{figure}

% =====================================================================
\section{Root-cause analysis: why the warm start diverges}
% =====================================================================
A time-marching code \emph{should} relax any initial condition to the same steady state,
so a divergence is not inevitable --- it points to a specific, identifiable mechanism. We
instrumented the base cells (an env-gated diagnostic, \texttt{EXHALE\_DIAG\_BASE=1}, dumping
$r,n,v,T$, heat, cool for the lowest cells over the first 500 steps; run with 60 OpenMP
threads) and traced the exact onset (Fig.~\ref{fig:basebreak}). The causal chain is:

\begin{enumerate}\itemsep3pt
\item \textbf{Inverted base pressure gradient (built into the IC$\,$+$\,$BC pairing).} The
  lower boundary condition (\texttt{Apply\_BC.f90}) pins the ghost cells to the \emph{cold}
  base state: $\rho_{\rm ghost}=\rho_{\rm bc}$, $p_{\rm ghost}=1+\delta p_{\rm bc}$ (hence
  $T_{\rm ghost}=T_0=1320$\,K), and $v_{\rm ghost}=\max(v_1,0)$. This is exactly consistent
  with a cold hydrostatic IC. But the warm IC places the first \emph{interior} cells at the
  warm pressure $p\!\approx\!1.35$ ($T\!\approx\!1840$\,K, $\sim\!1.4\times$ the cold value at
  the same density). So $p_{\rm interior}>p_{\rm ghost}$ --- an \emph{inverted} pressure
  gradient at the base, whose force points \emph{inward} and adds to gravity.

\item \textbf{The base flips to inflow within 3 steps} (Fig.~\ref{fig:basebreak}a). The base
  velocity reverses from the IC outflow ($+0.4$\,km\,s$^{-1}$) to inflow
  ($-0.09$\,km\,s$^{-1}$ at step 3) and then grows \emph{monotonically} more negative,
  reaching $-2.5$\,km\,s$^{-1}$ by step 480. The BC's $v_{\rm ghost}=\max(v_1,0)$ clamp acts
  as a \textbf{one-way valve}: it permits the inflow but supplies \emph{zero} outflow to
  relieve it, so the base never re-establishes the wind.

\item \textbf{The base compresses, it does not drain} (Fig.~\ref{fig:basebreak}b). The inflow
  piles gas up: $n_{\rm base}$ rises from $5.4\times10^{12}$ to $9.8\times10^{12}$ over 500
  steps. (This corrects the naive ``density-collapse'' guess --- the failure is compression,
  not evacuation.)

\item \textbf{A huge net heating with no advective escape} (Fig.~\ref{fig:basebreak}d). The
  base is photoionization-dominated: heat$/$cool $\approx 2000$. In a \emph{converged} wind
  this enormous net heating is carried off by the \emph{outflow} (advection $+$ $P\,\mathrm dV$
  expansion cooling). An \emph{inflowing}, compressing base has no such escape --- $P\,\mathrm
  dV$ now does positive (compressional) work and advection imports energy --- so the net heat
  accumulates.

\item \textbf{Slow $T$ climb, then runaway over the cooling peak} (Fig.~\ref{fig:basebreak}c).
  $T_{\rm base}$ rises steadily ($1675\to2710$\,K over 500 steps); it has \emph{not} yet
  reached the cooling peak at step 500 --- the full runaway comes much later
  ($T\to10^{40}\to$ NaN near step $30\,000$). The trigger is in the energy solver: in
  \texttt{energy\_semi\_implicit.f90} the implicit factor is
  $\mathrm dF/\mathrm dT = 1 + c_{\rm factor}\,\max(0,\,\mathrm dC/\mathrm dT)$, which is
  unconditionally stable \emph{only on the rising branch} of the cooling curve. Once a base
  cell is pushed past the cooling peak ($\sim$1.5--$2\times10^4$\,K for photoionized H/He,
  where the neutral fraction collapses and the cooling rate \emph{falls} with $T$, i.e.\
  $\mathrm dC/\mathrm dT\le0$), the $\max(0,\cdot)$ disengages the safeguard,
  $\mathrm dF/\mathrm dT\to1$, the update becomes \emph{fully explicit with no damping}, and
  the temperature runs away. \emph{(Since fixed: \texttt{energy\_semi\_implicit.f90} now uses
  $1 + c_{\rm factor}\,|\mathrm dC/\mathrm dT|$, i.e.\ recommendation~2 of
  \S\ref{sec:reco} below, so this branch of the mechanism no longer exists in the code.)}
\end{enumerate}

\noindent So the divergence is the product of two design assumptions, each individually
reasonable: a base BC \emph{hard-wired to the cold state}, and an energy solver whose
unconditional stability holds \emph{only on the rising branch} of the cooling curve. The
cold hydrostatic IC sits in the basin of attraction of both --- it matches the cold BC (no
inverted gradient, no base inflow) \emph{and} it approaches equilibrium from below (never
visiting the $\mathrm dC/\mathrm dT\le0$ branch). The warm Parker IC violates both at once.

\begin{figure}[ht]
\centering
\includegraphics[width=\linewidth]{../benchmarks/ic_start_benchmark/base_breakdown.pdf}
\caption{Instrumented hot-Parker startup (base cell $j{=}1$, first 500 steps, 60 threads).
(a) the base velocity flips from the IC outflow to \emph{inflow} by step 3 and grows
monotonically more negative (the one-way-valve base BC cannot relieve it); (b) the base
\emph{compresses} (density piles up) rather than draining; (c) $T$ climbs steadily with no
advective escape; (d) the base carries a near-constant net photoheating heat$/$cool
$\approx2000$. The cell has not yet crossed the cooling peak at step 500; once it does, the
semi-implicit energy safeguard disengages and $T$ runs away to NaN ($\sim$step $30\,000$).}
\label{fig:basebreak}
\end{figure}

% =====================================================================
\section{Making the warm start converge (implemented fix)}
% =====================================================================
The root-cause analysis (\S6) named two candidate triggers --- the cold-pinned base BC and
the energy-solver branch. Acting on them showed that the \emph{decisive} one is neither: it
is the \textbf{IC density}.

\paragraph{First attempt (insufficient): cold-base window on the full Parker IC.} Forcing the
base layer cold and static with a spatial smoothstep window (cold below a transition radius
$r_{\rm tr}$, warm Parker above) made the IC \emph{start} correct (base $T=T_0$, $v\!\approx\!0$,
$p=1+\delta p_{\rm bc}$), but the run still diverged: the base inflow \emph{re-developed
dynamically}, growing monotonically to $-2.7\times10^5$\,cm\,s$^{-1}$ over the first 500 steps.
So the inverted IC gradient was not the sustaining cause.

\paragraph{The diagnostic that settled it.} Running the \emph{cold hydrostatic} IC with the
same base instrumentation revealed that it, too, breathes inward --- the base velocity is
negative in 494 of the first 500 steps. Base breathing is therefore \emph{normal}. The
difference is the dynamics: the cold-hydrostatic inflow is \textbf{self-limiting} (it peaks at
$-2.6\times10^4$\,cm\,s$^{-1}$ near step 240 and then \emph{recedes}), whereas the warm-Parker
inflow is \textbf{self-amplifying} ($\sim\!10\times$ larger, monotonic). The culprit is the
\emph{isothermal-Parker density}: it is globally inconsistent with ATES's non-isothermal
equilibrium --- the hot, low-density Parker column is not supported against gravity at the cold
base and drains onto it, feeding the inflow.

\paragraph{The fix that works: keep the hydrostatic density, warm only the overlay.} The
working warm-seed IC (now what \texttt{Hot Parker IC} builds) keeps the \textbf{stable cold
hydrostatic density} --- the profile the cold IC proves is balanced --- and overlays on it only
\emph{(i)} a warm temperature $+$ H/He ionization ramp (which shortcuts the stiff
thermal/ionization relaxation) and \emph{(ii)} the Parker \emph{velocity} as a transonic
head-start, both anchored to the cold static base by the smoothstep window ($x_i=0$ below $r=1$,
$\to1$ at $r\ge$\,\texttt{hp\_base\_rtr}). The Parker solve is used for the velocity only; its
density is discarded. With this, the base inflow becomes self-limiting again (peak
$-2.8\times10^4$, then recedes --- indistinguishable from the cold-hydrostatic base), and the
run converges with no NaN.

\paragraph{Result: it converges to the same state, but does not speed up.} A clean head-to-head
(same planet, 60 OpenMP threads, both to convergence) gives Table~\ref{tab:warmfix}: the
warm-seed IC converges to the \emph{identical} steady state ($\log_{10}\dot M=11.52$ for both)
in essentially the \emph{same} cost ($93\,068$ vs $95\,274$ steps, $-2.3\%$; wall-clock
$+0.5\%$ --- within noise). Figure~\ref{fig:warmfix}a shows why: although the warm seed's early
transient is markedly \emph{gentler} (its $du$ peaks near $10^3$ versus $\sim\!10^6$ for the cold
start), the two convergence histories merge after $\sim\!15\,000$ steps and the long
slowly-decaying tail --- which sets the total step count --- is identical. The stopping
criterion inherited from ATES is the \textbf{hydrodynamic momentum-steadiness} test
($du<du_{\rm th}$, i.e.\ $\rho v r^2=$const; the runs of this benchmark used the loosened
$du_{\rm th}=2\times10^{-2}$ of the time, and the default is now $10^{-3}$), and the warm seed
accelerates the \emph{thermal} relaxation, not that
hydrodynamic one. So for this (momentum-gated) planet the warm start is a correctness-neutral
option, not a performance win.

\begin{table}[ht]
\centering
\caption{Cold hydrostatic vs.\ the fixed warm-seed IC (\texttt{Hot Parker IC}), both run to
convergence at 60 OpenMP threads. Same converged $\dot M$; no meaningful speedup.}
\label{tab:warmfix}
\small
\begin{tabular}{@{}lrrll@{}}
\toprule
IC & Steps & Wall [s] & Converged & $\log_{10}\dot M$ \\
\midrule
cold hydrostatic     & 95\,274 & 1851.1 & yes ($du{=}2.0\mathrm{e}{-}2$) & 11.52 \\
warm-seed (fixed)    & 93\,068 & 1860.9 & yes, no NaN                     & 11.52 \\
\bottomrule
\end{tabular}
\end{table}

\begin{figure}[ht]
\centering
\includegraphics[width=\linewidth]{../benchmarks/ic_start_benchmark/warmseed_fix.pdf}
\caption{The fixed warm-seed IC. (a) $du$ convergence history: the warm seed's early transient
is gentler (peak $\sim\!10^3$ vs $\sim\!10^6$), but the histories merge and the step-count-setting
tail is identical --- convergence is momentum-gated. (b,c) the warm-seed IC (dashed: already
warm/ionized, with a Parker velocity head-start) relaxes to the correct non-isothermal converged
profile (solid).}
\label{fig:warmfix}
\end{figure}

% =====================================================================
\section{Discussion and recommendation}\label{sec:reco}
% =====================================================================
\paragraph{Does a Parker-wind start speed up ATES? No --- and that is consistent with
how ATES converges.} ATES is a \emph{time-marching} code: it relaxes any initial
condition to the same steady state, and the IC only sets the cost of getting there. The
two relaxation channels (\S1) respond differently to a Parker seed:
\begin{itemize}\itemsep2pt
\item The \textbf{hydrodynamic} channel \emph{would} benefit from a transonic seed --- but
  only if a consistent transonic IC can be imposed at the base. For this planet the cold
  Parker has no interior sonic point, and the warm Parker conflicts with the base BC.
\item The \textbf{thermal/ionization} channel is the stiff, dominant cost. A warm/ionized
  seed could shortcut it --- but only if it does not break the base. It broke the base.
\end{itemize}

\paragraph{Why the warm seed was fragile --- and what actually fixed it.} The original
divergence (\S6) was real and proceeded exactly as traced (inverted base gradient $\to$ base
inflow $\to$ no advective escape $\to$ slow heating $\to$ runaway past the cooling peak). But
\S7 showed that the \emph{sustaining} cause was deeper than the inverted IC gradient: even with
the base started cold, the inflow re-developed, because the isothermal-Parker \emph{density} is
globally inconsistent with the non-isothermal equilibrium and drains onto the base. The
decisive fix was therefore to discard the Parker density and keep the stable cold hydrostatic
density, warming only the temperature/ionization/velocity overlay. That converges to the
correct steady state (Table~\ref{tab:warmfix}).

\paragraph{Recommendation.} Keep the \textbf{cold hydrostatic IC as the default}: even the
\emph{fixed} warm seed gives no speedup for this momentum-gated planet (\S7), so there is no
performance reason to switch. The \texttt{Hot Parker IC} keyword now builds the
\emph{validated} warm-seed IC (stable hydrostatic density $+$ warm overlay, cold-base smoothstep
window \texttt{hp\_base\_rtr}); it is a correct, converging option rather than an experimental
one, and may help on a problem whose convergence is \emph{thermally} gated rather than
momentum-gated (e.g.\ a denser, more strongly-cooled base where the thermal relaxation, not the
wind, sets the step count). Two further, independent improvements remain available:
\begin{enumerate}\itemsep2pt
\item \textbf{A characteristic / pressure-gradient base BC} (instead of the hard cold-pinned
  $\rho,p$ with $v=\max(v_1,0)$) would let the base breathe physically and widen the basin of
  attraction for non-hydrostatic starts generally.
\item \textbf{Make the energy solver unconditionally stable on \emph{both} branches}
  \emph{(done: the code now takes $|\mathrm dC/\mathrm dT|$)} --- use a
  true backward-Euler step on the full cooling (or take $|\mathrm dC/\mathrm dT|$ rather than
  $\max(0,\mathrm dC/\mathrm dT)$ in \texttt{energy\_semi\_implicit.f90}), so a cell pushed past
  the cooling peak during \emph{any} transient cannot run away. This is the same spirit as the
  option-(c) breathing-base guard already in the \texttt{\_adv} post-processor, applied to the
  main integration. (Not required for the warm-seed IC above, which no longer reaches the peak,
  but a worthwhile robustness margin.)
\end{enumerate}

The instrumentation used here is retained (env-gated \texttt{EXHALE\_DIAG\_BASE=1};
zero effect when unset) so any future base-stability work can re-trace the startup directly.

\paragraph{Contrast with relaxation BVP codes.} This negative result actually highlights an
ATES strength. Steady-state relaxation codes such as \texttt{wind-ae} (Numerical Recipes
\texttt{solvde} two-point BVP) \emph{require} a transonic initial guess close to the target
solution --- their own documentation warns the relaxation ``is incredibly sensitive to
being given a good initial guess \ldots\ if the guess is too far from the goal, the
relaxation method will not converge,'' and they ship parameter ``rampers'' to walk a known
solution to a new one. ATES needs no such guess: it converges from a generic cold
hydrostatic atmosphere precisely because it marches in time. The cost is that the stiff
thermal relaxation is slow; the benefit is robustness. The Parker-start experiment shows
that trying to buy the BVP-code's head start, without their careful base treatment, trades
that robustness away.

\end{document}
